%% ============================================================
%% LearnFlow -- Architecture Decision Records (ADR-001 bis ADR-007)
%% Compile: pdflatex LearnFlow_ADRs.tex  (zweimal fuer korrekte Referenzen)
%% ============================================================
\documentclass[a4paper, 11pt, oneside]{report}

%% --- Encoding & Sprache
\usepackage[T1]{fontenc}
\usepackage[utf8]{inputenc}
\usepackage[ngerman]{babel}

%% --- Schriften & Typographie
\usepackage{lmodern}
\usepackage{microtype}

%% --- Seitengeometrie
\usepackage[a4paper, top=25mm, bottom=25mm, left=28mm, right=22mm,
            headheight=15pt]{geometry}

%% --- Farben
\usepackage{xcolor}
\definecolor{adrblue}{RGB}{13,71,161}
\definecolor{adrlightblue}{RGB}{232,245,255}
\definecolor{adrline}{RGB}{100,149,237}
\definecolor{codebg}{RGB}{250,250,250}
\definecolor{codeframe}{RGB}{205,205,205}
\definecolor{kw}{RGB}{0,100,180}
\definecolor{str}{RGB}{163,21,21}
\definecolor{cmt}{RGB}{80,120,70}
\definecolor{statusproposed}{RGB}{240,150,0}
\definecolor{statusaccepted}{RGB}{40,140,60}
\definecolor{notebg}{RGB}{255,249,224}
\definecolor{noteframe}{RGB}{200,150,0}
\definecolor{warnbg}{RGB}{255,242,242}
\definecolor{warnframe}{RGB}{190,50,50}
\definecolor{tabhead}{RGB}{224,238,255}
\definecolor{depbg}{RGB}{255,244,244}

%% --- Mathe & Symbole
\usepackage{amsmath}
\usepackage{amssymb}
\usepackage{xspace}

%% --- Tabellen
\usepackage{booktabs}
\usepackage{tabularx}
\usepackage{array}
\usepackage{colortbl}
\usepackage{multirow}
\usepackage{longtable}

%% --- Listings
\usepackage{listings}

\lstdefinelanguage{TypeScript}{
  morekeywords={type,interface,class,const,let,var,async,await,return,if,
    else,new,implements,constructor,function,export,default,import,from,as,
    extends,private,public,protected,readonly,abstract,static,Promise,void,
    string,number,boolean,null,undefined,true,false,switch,case,break,for,
    of,in,throw,try,catch,finally,typeof,keyof,never,any,unknown,enum,
    declare,namespace,React,ReactNode,NextFunction,Request,Response,Map,
    Set,Record,Partial,Required,Pick,Omit},
  morestring=[b]",
  morestring=[b]',
  morestring=[b]`,
  morecomment=[l]{//},
  morecomment=[s]{/*}{*/},
  sensitive=true
}

\lstdefinelanguage{LearnSQL}{
  morekeywords={CREATE,TABLE,INSERT,INTO,SELECT,UPDATE,DELETE,FROM,WHERE,
    JOIN,LEFT,ON,AND,OR,NOT,NULL,IS,INDEX,PRIMARY,KEY,UNIQUE,FOREIGN,
    REFERENCES,DEFAULT,CONSTRAINT,VALUES,SET,ALTER,DROP,CASCADE,EXTENSION,
    IF,EXISTS,ORDER,BY,LIMIT,GROUP,DISTINCT,AS,USING,BIGSERIAL,UUID,TEXT,
    INTEGER,BOOLEAN,TIMESTAMPTZ,VECTOR,CHECK,IN,HNSW,SERIAL,WITH,REPLACE,
    RETURNING},
  morestring=[b]',
  morecomment=[l]{--},
  sensitive=false
}

\lstset{
  basicstyle      = \small\ttfamily,
  backgroundcolor = \color{codebg},
  frame           = single,
  rulecolor       = \color{codeframe},
  numbers         = left,
  numberstyle     = \tiny\color{gray!55},
  numbersep       = 10pt,
  xleftmargin     = 18pt,
  xrightmargin    = 4pt,
  breaklines      = true,
  breakatwhitespace = true,
  showstringspaces  = false,
  tabsize         = 2,
  keywordstyle    = \color{kw}\bfseries,
  stringstyle     = \color{str},
  commentstyle    = \color{cmt}\itshape,
  captionpos      = b,
  aboveskip       = \smallskipamount,
  belowskip       = \smallskipamount,
  extendedchars   = true,
  literate        =
    {ä}{{\"{a}}}1  {ö}{{\"{o}}}1  {ü}{{\"{u}}}1
    {Ä}{{\"{A}}}1  {Ö}{{\"{O}}}1  {Ü}{{\"{U}}}1
    {ß}{{\ss}}1    {é}{{\'e}}1    {è}{{\`e}}1
    {à}{{\`a}}1    {â}{{\^a}}1    {ê}{{\^e}}1
    {î}{{\^i}}1    {ô}{{\^o}}1    {û}{{\^u}}1,
}

\lstdefinestyle{diagram}{
  basicstyle      = \footnotesize\ttfamily,
  backgroundcolor = \color{gray!7},
  frame           = single,
  rulecolor       = \color{gray!35},
  numbers         = none,
  xleftmargin     = 6pt, xrightmargin = 6pt,
  breaklines      = false,
  keywordstyle    = {},
  stringstyle     = {},
  commentstyle    = {},
  identifierstyle = {},
  showstringspaces = false,
}

%% --- tcolorbox
\usepackage[breakable,skins]{tcolorbox}

%% --- Kopf-/Fusszeilen
\usepackage{fancyhdr}
\pagestyle{fancy}
\fancyhf{}
\fancyhead[L]{\small\itshape\color{gray!65}LearnFlow --- Architecture Decision Records}
\fancyhead[R]{\small\itshape\color{gray!65}\leftmark}
\fancyfoot[C]{\small\color{gray!65}\thepage}
\renewcommand{\headrulewidth}{0.35pt}

%% --- Kapitel- und Abschnittsformatierung
\usepackage{titlesec}

\titleformat{\chapter}[display]
  {\normalfont\sffamily\color{adrblue}}
  {\large\bfseries\color{adrblue!75}%
   ADR-\ifnum\value{chapter}<10 00\fi\arabic{chapter}}
  {3pt}
  {\LARGE\bfseries\color{adrblue}}
  [\vspace{2pt}{\color{adrline}\rule{\linewidth}{1.2pt}}\vspace{-4pt}]

\titlespacing*{\chapter}{0pt}{8pt}{14pt}

\titleformat{\section}
  {\normalfont\large\bfseries\color{adrblue!85!black}}
  {\thesection}{0.6em}{}
\titlespacing*{\section}{0pt}{11pt}{5pt}

\titleformat{\subsection}
  {\normalfont\normalsize\bfseries\color{black!80}}
  {\thesubsection}{0.5em}{}
\titlespacing*{\subsection}{0pt}{7pt}{3pt}

%% --- Hyperref (zuletzt)
\usepackage[colorlinks=true, linkcolor=adrblue, urlcolor=adrblue!80!black,
  pdftitle={LearnFlow -- Architecture Decision Records},
  pdfauthor={Christoph A. Amstutz, MD-PhD}]{hyperref}

%% --- Aufzaehlungen
\usepackage{enumitem}
\setlist{noitemsep, topsep=3pt}

%% --- Eigene Befehle
\newcommand{\code}[1]{\texttt{\small #1}}
\newcommand{\ra}{$\rightarrow$\xspace}
\newcommand{\todo}{$\square$\;}

\newcommand{\statusProposed}{%
  \colorbox{statusproposed!28}{\small\bfseries\color{statusproposed!70!black}\;Proposed\;}}
\newcommand{\statusAccepted}{%
  \colorbox{statusaccepted!28}{\small\bfseries\color{statusaccepted!60!black}\;Accepted\;}}

\newcommand{\abgelehnt}{\colorbox{depbg}{\small\bfseries\color{warnframe}Abgelehnt}\;}
\newcommand{\zurueckgestellt}{\colorbox{notebg}{\small\bfseries\color{noteframe}Zur\"{u}ckgestellt}\;}
\newcommand{\empfohlen}{\colorbox{statusaccepted!18}{\small\bfseries\color{statusaccepted!60!black}Empfohlen}\;}

%% Spaltentypen
\newcolumntype{L}[1]{>{\raggedright\arraybackslash}p{#1}}
\newcolumntype{C}[1]{>{\centering\arraybackslash}p{#1}}

%% Absatzabstand
\setlength{\parskip}{4pt plus 1pt}
\setlength{\parindent}{0pt}

%% ============================================================
\begin{document}
%% ============================================================

%% --- Titelseite
\begin{titlepage}
\centering
\vspace*{2.5cm}
{\color{adrblue}\rule{\linewidth}{2pt}}
\vspace{0.9cm}

{\Huge\bfseries\sffamily\color{adrblue} LearnFlow}\\[0.4cm]
{\LARGE\sffamily Architecture Decision Records}\\[0.25cm]
{\large ADR-001 bis ADR-007}

\vspace{0.8cm}
{\color{adrblue}\rule{\linewidth}{0.5pt}}
\vspace{0.8cm}

\begin{tabular}{@{}ll@{}}
\textbf{Version}  & 1.0 \\[4pt]
\textbf{Datum}    & 2026-05-27 \\[4pt]
\textbf{Autor}    & Christoph A. Amstutz, MD-PhD \\[4pt]
\textbf{Status}   & Alle ADRs: \statusProposed \\[4pt]
\textbf{Basis}    & \textit{LearnFlow\_Requirements\_v1.md} \\
\end{tabular}

\vspace{1.6cm}
\begin{tcolorbox}[
  enhanced, colback=adrlightblue, colframe=adrblue,
  boxrule=1pt, arc=4pt, width=0.85\linewidth]
\small Dieses Dokument enthält die Architecture Decision Records (ADRs)
für das MVP der LearnFlow-Lernplattform.
ADRs dokumentieren signifikante Architekturentscheidungen,
die betrachteten Alternativen sowie die Begründung für den gewählten Ansatz.
Alle ADRs haben den Status \textbf{Proposed} und sind zur Review freigegeben.
\end{tcolorbox}

\vfill
{\color{adrblue!40}\rule{\linewidth}{0.5pt}}\\[4pt]
{\small\color{gray!65} Erstellt mit Claude --- Anthropic $\cdot$ 2026-05-27}
\end{titlepage}

%% --- Inhaltsverzeichnis
\tableofcontents
\clearpage

%%============================================================
\chapter{Konfidenz-Scoring: Strategy Pattern mit austauschbarem Scorer}
\label{chap:adr001}
%%============================================================

\begin{tcolorbox}[enhanced, colback=adrlightblue, colframe=adrblue,
  boxrule=0.9pt, arc=3pt, left=8pt, right=8pt, top=6pt, bottom=6pt]
\begin{tabular}{@{}L{3.2cm}L{11.5cm}@{}}
\textbf{Status}     & \statusProposed \\[3pt]
\textbf{Datum}      & 2026-05-27 \\[3pt]
\textbf{Autor}      & Christoph A. Amstutz, MD-PhD \\[3pt]
\textbf{Bezug}      & US-02, QA-06 (Testability), QA-04 (Reliability) \\[3pt]
\textbf{Massnahme}  & M2 \\
\end{tabular}
\end{tcolorbox}

\section{Kontext}

US-02 schreibt eine dreistufige Unterdrückungslogik für KI-Antworten vor.
Der Entscheid vom 2026-05-20 legt die Vorrang-Reihenfolge fest:

\begin{enumerate}
  \item \textbf{Quellenprüfung} --- kein Retrieval-Treffer \ra Antwort wird nicht angezeigt
  \item \textbf{Konfidenz-Score} --- Score unter Schwellenwert \ra ,,Ich weiss es nicht``
  \item \textbf{Self-Check-Anteil} --- unter 80\,\% \ra ,,Eingeschränkt belegt``;
        unter 50\,\% \ra unterdrückt
\end{enumerate}

Risiko~2 aus dem Requirements-Dokument benennt das zentrale offene Problem:
\textit{wie} der Self-Check-Anteil konkret berechnet wird, ist nicht definiert.
Die zwei diskutierten Mechanismen --- semantische Ähnlichkeit zwischen
Antwort-Sätzen und Retrieval-Chunks vs.\ LLM-Selbstevaluation --- haben
unterschiedliche Qualitäts- und Kostenprofile. US-11 fordert zusätzlich,
dass Schwellenwerte ohne Code-Deployment änderbar sind.

\section{Entscheidung}

Der Konfidenz-Scoring-Mechanismus wird als \textbf{Strategy Pattern} implementiert:

\begin{itemize}
  \item Ein \code{ConfidenceScorer}-Interface kapselt ausschliesslich die
        Berechnung des Self-Check-Anteils.
  \item Ein \code{ConfidenceEvaluator} enthält die unveränderliche Stufenlogik
        und nimmt \code{ConfidenceScorer} sowie \code{ConfigService} als
        injizierte Abhängigkeiten entgegen.
  \item Schwellenwerte werden ausnahmslos über den \code{ConfigService} gelesen ---
        nie als Konstanten im Code.
\end{itemize}

\subsection{Komponentenstruktur}

\begin{lstlisting}[style=diagram]
ConfidenceEvaluator
  +-- ConfidenceScorer (Interface)
  |     +-- SemanticSimilarityScorer   [MVP]
  |     +-- LLMSelfEvaluationScorer    [Post-MVP]
  |     '-- FixedRatioScorer           [Tests]
  '-- ConfigService (Interface)
        +-- DBConfigService            [Produktion]
        '-- InMemoryConfigService      [Tests]
\end{lstlisting}

\subsection{Interface-Definition}

\begin{lstlisting}[language=TypeScript]
interface ConfidenceScorer {
  /** Gibt einen Wert zwischen 0.0 und 1.0 zurueck:
   *  Anteil der Antwort-Aussagen, die durch Retrieval-Chunks belegbar sind. */
  computeSelfCheckRatio(
    llmResponse:     string,
    retrievalChunks: Chunk[]
  ): Promise<number>
}
\end{lstlisting}

\subsection{Stufenlogik (ConfidenceEvaluator)}

\begin{lstlisting}[language=TypeScript]
type AnswerDecision =
  | { kind: 'NO_SOURCE'  }
  | { kind: 'DONT_KNOW'; hint: string }
  | { kind: 'SUPPRESSED' }
  | { kind: 'LIMITED'    }
  | { kind: 'OK'         }

class ConfidenceEvaluator {
  constructor(
    private scorer: ConfidenceScorer,
    private config: ConfigService
  ) {}

  async evaluate(hits: Chunk[], score: number, response: string)
      : Promise<AnswerDecision> {
    const t = await this.config.getThresholds()

    if (hits.length === 0)
      return { kind: 'NO_SOURCE' }                        // Stufe 1

    if (score < t.minConfidenceScore)
      return { kind: 'DONT_KNOW',
               hint: 'Versuche, einen konkreten Prozess zu nennen.' }  // Stufe 2

    const ratio = await this.scorer.computeSelfCheckRatio(response, hits)
    if (ratio < t.suppressThreshold) return { kind: 'SUPPRESSED' }     // Stufe 3a
    if (ratio < t.limitedThreshold)  return { kind: 'LIMITED'    }     // Stufe 3b
    return { kind: 'OK' }
  }
}
\end{lstlisting}

\subsection{MVP-Scorer: Semantische Ähnlichkeit}

Für das MVP wird \code{SemanticSimilarityScorer} implementiert:
jeder Satz der LLM-Antwort wird vektorisiert und mit der
Cosinus-Ähnlichkeit zu den Retrieval-Chunks verglichen.
Überschreitet ein Satz den Ähnlichkeitsschwellenwert, gilt er als
,,belegt``. Der Self-Check-Anteil ist der Quotient belegter Sätze
zu Gesamtsätzen.

\section{Begründung}

\paragraph{Warum Strategy Pattern statt Pure Function?}
Eine Pure Function \code{evaluateConfidence(input): AnswerDecision} wäre
ausreichend, wenn der Self-Check-Algorithmus feststünde. Da Risiko~2 explizit
offen ist und sich der Algorithmus nach dem Pilot voraussichtlich ändern wird,
würde eine Pure Function einen späteren Umbau erzwingen, der die gesamte
Stufenlogik tangiert. Das Strategy Pattern trennt die stabile Stufenlogik vom
instabilen Algorithmus. Kosten: eine Abstraktion und ca.\ 30 Minuten Mehraufwand.

\paragraph{Warum kein Chain-of-Responsibility?}
Drei fixe, geordnete Stufen rechtfertigen kein Pipeline-Pattern.
Der Mehraufwand für ein Context-Objekt und Handler-Verkettung übersteigt den
Nutzen für einen Anwendungsfall dieser Grösse.

\section{Betrachtete Alternativen}

\subsection{Alternative~1 -- Pure Function}

\begin{lstlisting}[language=TypeScript]
function evaluateConfidence(input: ScoringInput): AnswerDecision { ... }
\end{lstlisting}

\abgelehnt weil der Self-Check-Algorithmus (Parameter \code{selfCheckRatio})
von aussen übergeben werden müsste --- die Berechnung läge dann ungekapselt im
aufrufenden Code. Das verlagert das Problem, löst es nicht.

\subsection{Alternative~2 -- Chain of Responsibility}

Jede Stufe ist ein eigenständiger Handler; ein Context-Objekt wandert durch
die Kette.

\abgelehnt Drei fixe, sequenzielle Stufen mit eindeutiger Reihenfolge sind
kein Anwendungsfall für dieses Pattern. Erhöhter Boilerplate ohne Mehrwert.

\section{Konsequenzen}

\paragraph{Positiv}
\begin{itemize}
  \item Self-Check-Algorithmus austauschbar ohne Änderung an der Stufenlogik
  \item \code{ConfidenceEvaluator} mit \code{FixedRatioScorer} und
        \code{InMemoryConfigService} in unter 5 Unit-Tests vollständig abgedeckt
  \item Schwellenwerte nie als Magic Numbers im Code --- alle Änderungen laufen
        über \code{ConfigService} (\ra ADR-003)
  \item Post-MVP ML-basierter Scorer als neue Implementierung einhängbar
\end{itemize}

\paragraph{Negativ / Risiken}
\begin{itemize}
  \item Zwei Interfaces statt einer Funktion: minimaler Erklärungsaufwand
        beim Onboarding
  \item \code{SemanticSimilarityScorer} benötigt Zugriff auf den
        Embedding-Adapter (\ra Abhängigkeit auf LLM-Infrastruktur auch
        im Scoring-Pfad)
  \item Qualität des MVP-Scorers erst durch RAG-Evaluierungs-Harness
        validierbar --- kein theoretischer Nachweis möglich
\end{itemize}

\section{Abhängigkeiten}

\begin{table}[htb]
\small\centering
\begin{tabularx}{\linewidth}{L{4cm} L{3cm} X}
\toprule
\rowcolor{tabhead}
\textbf{Abhängigkeit} & \textbf{Typ} & \textbf{Hinweis} \\
\midrule
ADR-003 (\code{ConfigService})  & Voraussetzung & Schwellenwerte kommen aus \code{ConfigService} \\
ADR-002 (RAG-Pipeline)          & Verwendend    & \code{ConfidenceEvaluator} wird in der Pipeline aufgerufen \\
QA-06 Testability (M3)          & Validierung   & RAG-Evaluierungs-Harness prüft die 90\,\%-,,Weiss-ich-nicht``-Rate \\
Risiko~2                        & Offen         & Self-Check-Algorithmus muss vor Sprint~1 als Spike validiert werden \\
\bottomrule
\end{tabularx}
\end{table}

\section{Offen / Nächste Schritte}

\begin{itemize}[label=\todo]
  \item Spike: \code{SemanticSimilarityScorer} gegen Evaluierungs-Dataset
        ausführen; Qualitätsschwelle festlegen (Go/No-Go für Risiko 1+2)
  \item Konkrete Ähnlichkeitsschwelle für ,,Satz gilt als belegt`` definieren
        (Cosinus-Score $> X$)
  \item Entscheiden, ob \code{selfCheckRatio} im Response-Objekt für Lara
        sichtbar ist oder nur als interne Metainfo bleibt
\end{itemize}

%%============================================================
\chapter{Fail-Safe RAG-Pipeline: Discriminated Union als Ergebnistyp}
\label{chap:adr002}
%%============================================================

\begin{tcolorbox}[enhanced, colback=adrlightblue, colframe=adrblue,
  boxrule=0.9pt, arc=3pt, left=8pt, right=8pt, top=6pt, bottom=6pt]
\begin{tabular}{@{}L{3.2cm}L{11.5cm}@{}}
\textbf{Status}     & \statusProposed \\[3pt]
\textbf{Datum}      & 2026-05-27 \\[3pt]
\textbf{Autor}      & Christoph A. Amstutz, MD-PhD \\[3pt]
\textbf{Bezug}      & US-01, US-02, QA-04 (Reliability), QA-06 (Testability) \\[3pt]
\textbf{Massnahme}  & M5 \\
\end{tabular}
\end{tcolorbox}

\section{Kontext}

Die RAG-Pipeline ist der kritische Pfad von LearnFlow:

\begin{lstlisting}[style=diagram]
Eingabe-Validierung -> Retrieval -> Konfidenz-Check
                    -> LLM-Call -> Self-Check -> Antwort-Assembly
\end{lstlisting}

An jedem Schritt kann ein anderer Fehler auftreten, der eine andere Reaktion erfordert:

\begin{table}[htb]
\small\centering
\begin{tabularx}{\linewidth}{L{4cm} X X}
\toprule
\rowcolor{tabhead}
\textbf{Fehlerzustand} & \textbf{Korrekte Reaktion} & \textbf{Falsche Reaktion} \\
\midrule
Kein Retrieval-Treffer   & ,,Keine Quelle`` anzeigen           & Antwort ohne Quelle generieren \\
Konfidenz zu tief        & ,,Ich weiss es nicht`` + Hinweis    & Antwort mit tiefer Konfidenz ausgeben \\
LLM-API nicht erreichbar & Klare Fehlermeldung                 & Leere oder halluzinierte Antwort \\
LLM-Timeout              & Fehlermeldung nach 1 Retry          & Unbegrenztes Warten \\
Self-Check $< 50\,\%$    & Antwort vollständig unterdrücken    & Unmarkierte Antwort ausgeben \\
\bottomrule
\end{tabularx}
\end{table}

Das Requirements-Dokument formuliert: \textit{,,Eine falsch dargestellte Antwort
ist schlimmer als gar keine.``} Das Kernproblem mit Exception-Hierarchien:
ein vergessener \code{catch}-Branch gibt implizit \code{undefined} oder einen
Teilzustand zurück --- der exakte Mechanismus, der halluzinierte Antworten
bei Ausfällen ermöglicht.

\section{Entscheidung}

Die RAG-Pipeline gibt \textbf{immer} einen explizit typisierten
\code{RAGResult}-Wert zurück. Der Rückgabetyp ist ein
\textbf{Discriminated Union}: eine abgeschlossene Menge benannter
Varianten, die alle möglichen Pipeline-Ausgaben beschreiben.

\subsection{Ergebnistyp}

\begin{lstlisting}[language=TypeScript]
type RAGResult =
  | { kind: 'answer'; content: string; sources: Source[];
      confidence: number; label: 'ok' | 'limited' }
  | { kind: 'dont_know';    hint: string    }
  | { kind: 'no_source'                     }
  | { kind: 'suppressed'                    }
  | { kind: 'unavailable';  message: string }
  | { kind: 'timeout'                       }
\end{lstlisting}

Jede Variante entspricht direkt einem UI-Zustand. Es gibt keine Variante,
die eine halluzinierte oder ungeprüfte Antwort transportiert.

\subsection{Pipeline-Implementierung}

\begin{lstlisting}[language=TypeScript]
async function runRAGPipeline(question: string): Promise<RAGResult> {
  const hits = await retriever.search(question) // wirft nie -- [] = kein Treffer

  const score    = computeRetrievalScore(hits)
  const decision = await evaluator.evaluate(hits, score, '')
  if (decision.kind === 'NO_SOURCE') return { kind: 'no_source'  }
  if (decision.kind === 'DONT_KNOW') return { kind: 'dont_know', hint: decision.hint }

  const llmResult = await llmAdapter.complete(buildPrompt(question, hits))
  if (!llmResult.ok)
    return llmResult.error.kind === 'timeout'
      ? { kind: 'timeout' }
      : { kind: 'unavailable', message: 'KI-Service nicht erreichbar.' }

  const final = await evaluator.evaluate(hits, score, llmResult.value.text)
  switch (final.kind) {
    case 'SUPPRESSED': return { kind: 'suppressed' }
    case 'LIMITED':    return { kind: 'answer', content: llmResult.value.text,
                                sources: extractSources(hits),
                                confidence: score, label: 'limited' }
    case 'OK':         return { kind: 'answer', content: llmResult.value.text,
                                sources: extractSources(hits),
                                confidence: score, label: 'ok' }
    default:           return { kind: 'unavailable', message: 'Unerwarteter Zustand.' }
  }
}
\end{lstlisting}

\subsection{Exhaustiveness im UI-Controller}

\begin{lstlisting}[language=TypeScript]
// Fehlender Case => TypeScript-Compilerfehler -- kein Runtime-Problem
function renderResult(result: RAGResult): ReactNode {
  switch (result.kind) {
    case 'answer':      return result.label === 'limited'
                          ? <LimitedAnswerCard {...result} />
                          : <AnswerCard {...result} />
    case 'dont_know':   return <DontKnowCard hint={result.hint} />
    case 'no_source':
    case 'suppressed':  return <NoSourceCard />
    case 'unavailable': return <ErrorCard message={result.message} />
    case 'timeout':     return <ErrorCard message="Zeitüberschreitung." />
  }
}
\end{lstlisting}

\subsection{LLM-Adapter-Interface mit Result-Typ}

\begin{lstlisting}[language=TypeScript]
type LLMError =
  | { kind: 'timeout'   }
  | { kind: 'api_error'; statusCode: number; detail: string }
  | { kind: 'rate_limit' }

type Result<T, E> = { ok: true; value: T } | { ok: false; error: E }

interface LLMAdapter {
  /** Gibt immer ein Result zurueck -- wirft nie. */
  complete(prompt: string): Promise<Result<LLMResponse, LLMError>>
  stream(prompt: string):   AsyncIterable<string>
}
\end{lstlisting}

\subsection{Flussdiagramm}

\begin{lstlisting}[style=diagram]
question
    |
    v
retriever.search()
    +-(leer)-----------------------------------------> { kind: 'no_source' }
    | hits
    v
evaluator.evaluate() [pre-LLM]
    +-(NO_SOURCE)-----------------------------------> { kind: 'no_source' }
    +-(DONT_KNOW)-----------------------------------> { kind: 'dont_know', hint }
    | pass
    v
llmAdapter.complete()
    +-(Err: timeout)--------------------------------> { kind: 'timeout' }
    +-(Err: api_error / rate_limit)-----------------> { kind: 'unavailable' }
    | Ok: llmResponse
    v
evaluator.evaluate() [post-LLM, mit Self-Check]
    +-(SUPPRESSED)----------------------------------> { kind: 'suppressed' }
    +-(LIMITED)-------------------------------------> { kind: 'answer', label:'limited' }
    '-(OK)------------------------------------------> { kind: 'answer', label:'ok' }
\end{lstlisting}

\section{Begründung}

\paragraph{Warum Discriminated Union statt Exception-Hierarchie?}
Alle sechs \code{RAGResult}-Varianten sind \textit{erwartete} Zustände ---
direkt aus den Akzeptanzkriterien von US-01 und US-02 ableitbar. Der
TypeScript-Compiler erzwingt den Exhaustiveness-Check im
\code{switch}-Statement. Ein vergessener \code{catch}-Branch kann
nicht zu einer halluzinierten Antwort führen.

\paragraph{Warum zwei \code{evaluator.evaluate()}-Aufrufe?}
Der erste Aufruf (pre-LLM) behandelt Stufe~1 und~2 --- ohne LLM-Call prüfbar,
spart API-Kosten bei frühen Exits. Der zweite Aufruf führt den
Self-Check (Stufe~3) durch, der die LLM-Antwort benötigt.

\section{Betrachtete Alternativen}

\subsection{Alternative~1 -- Exception-Hierarchie}

\begin{lstlisting}[language=TypeScript]
class RAGException       extends Error {}
class NoSourceException  extends RAGException {}
class LLMUnavailableException extends RAGException {}
\end{lstlisting}

\abgelehnt Kein Compiler-Exhaustiveness-Check. Vergessener \code{catch}-Branch
führt zu undefiniertem Verhalten. Kontrollfluss via Exceptions ist schwerer
nachzuvollziehen und zu testen.

\subsection{Alternative~2 -- Result-Monad (fp-ts / neverthrow)}

\abgelehnt für MVP\textit{:} Vollwertige Monad-Komposition erhöht die
Einstiegshürde für das Team erheblich. Der Mehrwert gegenüber einfachen
Discriminated Unions ist für diesen Anwendungsfall gering.
Für ein Python-Backend ist der native \code{match}-Ausdruck (ab 3.10)
die bessere Alternative.

\subsection{Alternative~3 -- Response-Objekt mit Status-Feld}

\begin{lstlisting}[language=TypeScript]
type RAGResponse = { status: 'ok' | 'error' | 'dont_know'; data?: any; error?: any }
\end{lstlisting}

\abgelehnt Optionale Felder (\code{?}) entfernen die Compiler-Garantie.
\code{result.data} könnte \code{undefined} sein, auch wenn
\code{status === 'ok'}.

\section{Konsequenzen}

\paragraph{Positiv}
\begin{itemize}
  \item Vollständige Abdeckung aller Fehlerpfade ist vom Compiler erzwingbar
  \item Jede \code{RAGResult}-Variante isoliert und ohne Netzwerk testbar
  \item UI-Komponenten direkt auf Varianten gemappt --- keine implizite
        Null-Prüfung nötig
  \item Neue Fehlerzustände können als neue Variante ergänzt werden;
        der Compiler zeigt sofort alle unbehandelten Stellen
\end{itemize}

\paragraph{Negativ / Risiken}
\begin{itemize}
  \item In Python nicht nativ (Klassen-Hierarchie oder \code{dataclass}-Varianten
        nötig); Exhaustiveness-Check muss manuell diszipliniert werden
  \item Die \code{Result<T, E>}-Abstraktion am LLM-Adapter ist für Entwickler
        ungewohnt, die ausschliesslich exception-basiert gearbeitet haben
  \item Jede neue \code{RAGResult}-Variante erfordert Anpassung im UI-Controller
\end{itemize}

\section{Abhängigkeiten}

\begin{table}[htb]
\small\centering
\begin{tabularx}{\linewidth}{L{4cm} L{3cm} X}
\toprule
\rowcolor{tabhead}
\textbf{Abhängigkeit} & \textbf{Typ} & \textbf{Hinweis} \\
\midrule
ADR-001 (\code{ConfidenceEvaluator}) & Verwendet   & Pipeline ruft \code{evaluator.evaluate()} zweimal auf \\
ADR-003 (\code{ConfigService})       & Indirekt    & Über ADR-001 --- Schwellenwerte aus Config \\
US-01 AC: Fehlermeldung bei Ausfall  & Validierung & \code{unavailable}-Variante deckt dieses Kriterium ab \\
US-02 AC: Unterdrückung $< 50\,\%$  & Validierung & \code{suppressed}-Variante deckt dieses Kriterium ab \\
QA-04 Circuit Breaker (M6)           & Vorgelagert & Circuit Breaker gibt \code{Err(rate\_limit)} zurück \\
\bottomrule
\end{tabularx}
\end{table}

\section{Offen / Nächste Schritte}

\begin{itemize}[label=\todo]
  \item Entscheiden: TypeScript oder Python für Backend? Bestimmt die Syntax des
        Discriminated Union
  \item \code{Source}-Typ definieren (Dokumenttitel, Abschnitt, Upload-Datum,
        Highlight-Offset für US-01)
  \item Timeout-Wert für LLM-Call konfigurierbar machen (\ra ADR-003)
  \item Retry-Logik im LLM-Adapter spezifizieren: 1 Retry mit wie viel Backoff?
\end{itemize}

%%============================================================
\chapter{Konfigurationsservice: PostgreSQL mit In-Process-Cache und Audit-Log}
\label{chap:adr003}
%%============================================================

\begin{tcolorbox}[enhanced, colback=adrlightblue, colframe=adrblue,
  boxrule=0.9pt, arc=3pt, left=8pt, right=8pt, top=6pt, bottom=6pt]
\begin{tabular}{@{}L{3.2cm}L{11.5cm}@{}}
\textbf{Status}     & \statusProposed \\[3pt]
\textbf{Datum}      & 2026-05-27 \\[3pt]
\textbf{Autor}      & Christoph A. Amstutz, MD-PhD \\[3pt]
\textbf{Bezug}      & US-02, US-06, US-11, QA-04 (Reliability), QA-05 (Maintainability) \\[3pt]
\textbf{Massnahme}  & M7 \\
\end{tabular}
\end{tcolorbox}

\section{Kontext}

Drei User Stories setzen voraus, dass Systemparameter zur Laufzeit änderbar
sind --- ohne Code-Deployment und ohne Applikations-Neustart:

\begin{table}[htb]
\small\centering
\begin{tabularx}{\linewidth}{L{5cm} L{2cm} L{3cm}}
\toprule
\rowcolor{tabhead}
\textbf{Parameter} & \textbf{Story} & \textbf{Standardwert} \\
\midrule
\code{min\_confidence\_score} & US-02 & TBD (via Spike) \\
\code{suppress\_threshold}    & US-02 & 0.50 \\
\code{limited\_threshold}     & US-02 & 0.80 \\
\code{stale\_days}            & US-06 & 90 \\
\bottomrule
\end{tabularx}
\end{table}

US-11 ergänzt: Konfigurationsänderungen werden mit Zeitstempel und
ausführender Person protokolliert. Änderungen müssen sofort wirken.
Das zentrale Problem mit ENV-Variablen oder Config-Files: sie erfordern
einen Neustart.

\section{Entscheidung}

Konfigurationsparameter werden in der \textbf{bestehenden PostgreSQL-Datenbank}
gespeichert. Ein zentraler \code{ConfigService} liest Werte über einen
\textbf{In-Process-Cache mit TTL und sofortiger Invalidierung beim Schreiben}.
Alle Schreiboperationen werden transaktionssicher in einem unveränderlichen
\textbf{Audit-Log} protokolliert.

\subsection{Datenbankschema}

\begin{lstlisting}[language=LearnSQL]
CREATE TABLE config_params (
  key        TEXT        PRIMARY KEY,
  value      TEXT        NOT NULL,
  updated_at TIMESTAMPTZ NOT NULL DEFAULT now()
);

INSERT INTO config_params (key, value) VALUES
  ('min_confidence_score', '0.65'),
  ('suppress_threshold',   '0.50'),
  ('limited_threshold',    '0.80'),
  ('stale_days',           '90');

-- Unveraenderliches Audit-Log (kein DELETE/UPDATE fuer die Applikation)
CREATE TABLE config_audit_log (
  id         BIGSERIAL   PRIMARY KEY,
  key        TEXT        NOT NULL,
  old_value  TEXT,                     -- NULL bei Erstanlage
  new_value  TEXT        NOT NULL,
  changed_by TEXT        NOT NULL,     -- Username (kein FK)
  changed_at TIMESTAMPTZ NOT NULL DEFAULT now()
);
\end{lstlisting}

\begin{table}[htb]
\small\centering
\begin{tabular}{@{}lll@{}}
\toprule
\rowcolor{tabhead}
\textbf{Tabelle} & \textbf{App-Rolle} & \textbf{Admin-UI-Rolle} \\
\midrule
\code{config\_params}    & SELECT, UPDATE         & SELECT, UPDATE \\
\code{config\_audit\_log} & SELECT, INSERT         & SELECT, INSERT \\
\bottomrule
\end{tabular}
\end{table}

Kein DELETE, kein TRUNCATE auf \code{config\_audit\_log} ---
weder für die App noch für die Admin-UI.

\subsection{Interface}

\begin{lstlisting}[language=TypeScript]
interface ConfigService {
  get(key: string):        Promise<string>
  getNumber(key: string):  Promise<number>
  getThresholds():         Promise<ConfidenceThresholds>
  set(key: string, value: string, changedBy: string): Promise<void>
}

type ConfidenceThresholds = {
  minConfidenceScore: number   // 'min_confidence_score'
  suppressThreshold:  number   // 'suppress_threshold'
  limitedThreshold:   number   // 'limited_threshold'
}
\end{lstlisting}

\subsection{Produktions-Implementierung (DBConfigService)}

\begin{lstlisting}[language=TypeScript]
class DBConfigService implements ConfigService {
  private cache = new Map<string, { value: string; expiresAt: number }>()
  private readonly TTL_MS = 30_000   // 30 Sekunden

  constructor(private db: Database) {}

  async get(key: string): Promise<string> {
    const cached = this.cache.get(key)
    if (cached && Date.now() < cached.expiresAt) return cached.value

    const row = await this.db.queryOne<{ value: string }>(
      'SELECT value FROM config_params WHERE key = $1', [key]
    )
    if (!row) throw new ConfigKeyNotFoundError(key)
    this.cache.set(key, { value: row.value, expiresAt: Date.now() + this.TTL_MS })
    return row.value
  }

  async set(key: string, newValue: string, changedBy: string): Promise<void> {
    const oldValue = await this.get(key).catch(() => null)
    await this.db.transaction(async (tx) => {
      await tx.execute(
        'UPDATE config_params SET value = $1, updated_at = now() WHERE key = $2',
        [newValue, key])
      await tx.execute(
        'INSERT INTO config_audit_log (key, old_value, new_value, changed_by) VALUES ($1,$2,$3,$4)',
        [key, oldValue, newValue, changedBy])
    })
    this.cache.delete(key)   // sofortige Invalidierung
  }
}
\end{lstlisting}

\subsection{Zeitachse einer Schwellenwert-Änderung}

\begin{lstlisting}[style=diagram]
t=0s  Admin setzt suppress_threshold = 0.60 (via Admin-UI)
t=0s  ConfigService.set() -> DB UPDATE + Audit-Log INSERT (1 Transaktion)
t=0s  cache.delete('suppress_threshold')
t~1s  Naechste Anfrage: cache miss -> DB-Lesezugriff -> neuer Wert aktiv
      => Keine Neustart, kein Deployment erforderlich

[Single-Instance: Aenderung wirkt innerhalb ~1 Sekunde]
[Multi-Instance (Post-MVP): bis zu 30s Inkonsistenz -> Redis-Upgrade]
\end{lstlisting}

\section{Begründung}

\paragraph{Warum nicht ENV-Variablen?}
ENV-Variablen werden beim Prozessstart eingelesen. Änderungen erfordern
einen Neustart --- widerspricht der expliziten MUST-Anforderung aus US-11.

\paragraph{Warum nicht Redis als primärer Config-Store?}
Im MVP ist eine einzige App-Instanz geplant (20 concurrent users). Redis als
neue Infrastrukturkomponente für ein 30-Sekunden-Propagationsproblem, das im
Pilotbetrieb nie auftreten wird, widerspricht dem Budget-Constraint.
Der Upgrade-Pfad ist dokumentiert und ohne Interface-Änderung nachrüstbar.

\paragraph{Warum sofortige Cache-Invalidierung statt reinem TTL?}
Rein TTL-basierter Cache bedeutet: direkt nach einer Admin-Änderung sieht
dieselbe Instanz für bis zu 30 Sekunden den alten Wert. Sofortige
Invalidierung bei \code{set()} kostet eine Zeile Code und eliminiert
das Problem vollständig für Single-Instance.

\paragraph{Warum kein FK von \code{changed\_by} auf \code{users}?}
Administratoren könnten gelöscht werden (Mitarbeitendenwechsel). Ein FK würde
den Löschvorgang blockieren oder das Audit-Log korrumpieren. Der Username
als Text überlebt User-Löschoperationen.

\section{Betrachtete Alternativen}

\subsection{Alternative~1 -- ENV-Variablen / Config-File}

\abgelehnt Änderungen erfordern Neustart. Verletzt MUST-Anforderung
aus US-11 direkt.

\subsection{Alternative~2 -- Redis mit Pub/Sub-Invalidierung}

\begin{lstlisting}[style=diagram]
Admin -> UPDATE config_params (PostgreSQL)
      -> PUBLISH config:invalidate (Redis)
Alle Instanzen -> SUBSCRIBE -> cache.delete(key)  [sofort]
\end{lstlisting}

\zurueckgestellt Korrekte Architektur für Multi-Instance-Betrieb.
Für MVP (Single-Instance) unverhältnismässig. Interface bleibt kompatibel ---
Upgrade ohne Business-Logik-Änderung möglich.
\textbf{Upgrade-Trigger:} mehr als eine App-Instanz im Betrieb.

\subsection{Alternative~3 -- Polling ohne Cache}

\abgelehnt Der \code{ConfidenceEvaluator} ruft \code{getThresholds()} bei
jeder Nutzeranfrage auf. Bei 20 concurrent users entstehen 60--100 DB-Queries
pro Sekunde allein für Config-Reads.

\section{Konsequenzen}

\paragraph{Positiv}
\begin{itemize}
  \item Keine neue Infrastrukturkomponente: Config in bestehender PostgreSQL-Instanz
  \item Audit-Log transaktionssicher: Änderung und Log gelingen oder scheitern gemeinsam
  \item \code{InMemoryConfigService} macht alle Komponenten ohne DB testbar
  \item Schwellenwerte nie als Magic Numbers im Code
\end{itemize}

\paragraph{Negativ / Risiken}
\begin{itemize}
  \item Bei zukünftigem Multi-Instance-Betrieb: bis zu 30~s Inkonsistenz
        (akzeptabel für MVP; bei Scale-Up adressieren)
  \item DB-Ausfall macht Config-Reads unmöglich: Fallback auf gecachte Werte
        (bis TTL abläuft) vorhanden, aber kein expliziter Kaltstart-Fallback
        definiert
\end{itemize}

\section{Abhängigkeiten}

\begin{table}[htb]
\small\centering
\begin{tabularx}{\linewidth}{L{4cm} L{3cm} X}
\toprule
\rowcolor{tabhead}
\textbf{Abhängigkeit} & \textbf{Typ} & \textbf{Hinweis} \\
\midrule
ADR-001 (\code{ConfidenceEvaluator}) & Konsument    & Liest \code{getThresholds()} bei jeder Anfrage \\
ADR-002 (RAG-Pipeline)               & Indirekt     & Über ADR-001 \\
US-11 Admin-Seite                    & Schreibweg   & Admin-UI ruft \code{ConfigService.set()} auf \\
QA-03 Security                       & Voraussetzung & Admin-Seite nur für Admin-Rolle erreichbar \\
PostgreSQL                           & Infrastruktur & Kein zusätzliches System erforderlich \\
\bottomrule
\end{tabularx}
\end{table}

\section{Offen / Nächste Schritte}

\begin{itemize}[label=\todo]
  \item Initiale Schwellenwerte für \code{min\_confidence\_score} nach
        RAG-Spike festlegen und in Seed-Migration eintragen
  \item DB-Migrationsscript für \code{config\_params} und
        \code{config\_audit\_log} erstellen (Teil von Sprint~0)
  \item Entscheiden: Admin-UI (US-11) zeigt Audit-Log als read-only Tabelle --- Scope für MVP?
  \item Kaltstart-Verhalten bei leerem Cache und DB-Ausfall definieren
        (Default-Werte im Code als letzter Fallback?)
\end{itemize}

%%============================================================
\chapter{LLM-Provider-Wahl: Azure OpenAI Service oder Anthropic Claude API}
\label{chap:adr004}
%%============================================================

\begin{tcolorbox}[enhanced, colback=adrlightblue, colframe=adrblue,
  boxrule=0.9pt, arc=3pt, left=8pt, right=8pt, top=6pt, bottom=6pt]
\begin{tabular}{@{}L{3.2cm}L{11.5cm}@{}}
\textbf{Status}      & \statusProposed \\[3pt]
\textbf{Datum}       & 2026-05-27 \\[3pt]
\textbf{Autor}       & Christoph A. Amstutz, MD-PhD \\[3pt]
\textbf{Bezug}       & US-01, US-02, QA-03 (Security), QA-04 (Reliability), Risiko~3 \\[3pt]
\textbf{Abhängigkeit}& ADR-002 (LLMAdapter-Interface), ADR-005 (RAG-Stack) \\
\end{tabular}
\end{tcolorbox}

\section{Kontext}

LearnFlow schickt interne Unternehmensdokumente als Kontext an ein LLM.
Das stellt die Provider-Wahl unter zwei Zwänge, die über reine
Qualitäts- und Kostenüberlegungen hinausgehen:

\textbf{Datenschutz / Datenresidenz:} Jeder Prompt enthält Ausschnitte aus
vertraulichen internen Dokumenten. Je nach Unternehmensstandort (CH/DE/AT)
kann die Übermittlung an einen US-amerikanischen API-Endpoint
compliance-relevant sein.

\textbf{API-Quota} (Risiko~3): Quota-Anfragen können mehrere Wochen dauern.
Die Provider-Entscheidung muss deshalb in Sprint~0 getroffen und die
Quota-Anfrage sofort gestellt werden.

Das LLM-Adapter-Interface aus ADR-002 stellt sicher, dass der Provider
später ohne Änderung an der Business-Logik ausgetauscht werden kann.

\subsection{Bewertungskriterien}

\begin{table}[htb]
\small\centering
\begin{tabularx}{\linewidth}{L{6cm} L{2.5cm} X}
\toprule
\rowcolor{tabhead}
\textbf{Kriterium} & \textbf{Gewicht} & \textbf{Begründung} \\
\midrule
Datenresidenz / DSGVO-Compliance        & Hoch   & Interne Dokumente im Prompt \\
RAG-Qualität (Long-Context, Präzision)  & Hoch   & Kernversprechen der Plattform \\
Quota-Verfügbarkeit für MVP             & Hoch   & Risiko~3 \\
Kosten (Token-Preis $\times$ Volumen)   & Mittel & 480h Budget, $\approx$20 User \\
Streaming-Support                       & Mittel & $p_{95} \leq 10\,\text{s}$ (QA-01) \\
Embedding-Verfügbarkeit beim Provider   & Mittel & Vereinfacht ADR-005 \\
\bottomrule
\end{tabularx}
\end{table}

\section{Entscheidung}

\textbf{Zweistufige Empfehlung abhängig von der Compliance-Anforderung:}

\subsection{Pfad~A -- Azure OpenAI Service \textit{(wenn EU/EWR-Datenresidenz gefordert)}}

Azure OpenAI bietet dieselben Modelle (GPT-4o, GPT-4o-mini) über
Microsoft-Infrastruktur in europäischen Rechenzentren an und stellt
einen EU-Datenverarbeitungsvertrag (DPA nach Art.~28 DSGVO) bereit.

\textbf{Empfohlene Modelle:}
\begin{itemize}
  \item \textbf{GPT-4o-mini} --- Hauptmodell: niedrige Latenz, günstiger Token-Preis
  \item \textbf{GPT-4o} --- Fallback für komplexe, mehrstufige Kontextfragen
\end{itemize}

\subsection{Pfad~B -- Anthropic Claude API \textit{(wenn US-Datenresidenz akzeptiert)}}

Claude 3.5 Haiku / Sonnet bieten hervorragende Qualität für
Retrieval-gestützte Antwortgenerierung mit präziser Quellenreferenzierung.

\textbf{Empfohlene Modelle:}
\begin{itemize}
  \item \textbf{Claude 3.5 Haiku} --- Hauptmodell: $p_{50} \approx 2{-}3\,\text{s}$,
        günstiger als Sonnet
  \item \textbf{Claude 3.5 Sonnet} --- Fallback für komplexe Fragen
\end{itemize}

\begin{tcolorbox}[enhanced, colback=warnbg, colframe=warnframe,
  boxrule=0.8pt, arc=2pt, left=8pt, right=8pt, fontupper=\small]
\textbf{Entscheidungsblockierer:} Die Compliance-Frage (Datenresidenz) muss
durch die IT-/Legal-Abteilung bis Ende Sprint~0 beantwortet werden.
Bis dahin bleibt dieser ADR \textbf{Proposed}. Die technische Implementierung
ist für beide Pfade identisch (LLM-Adapter-Interface, ADR-002).
\end{tcolorbox}

\subsection{Gemeinsame Konfiguration (beide Pfade)}

\begin{lstlisting}[language=TypeScript]
interface LLMAdapter {
  complete(prompt: string, options?: CompletionOptions)
      : Promise<Result<LLMResponse, LLMError>>
  stream(prompt: string, options?: CompletionOptions): AsyncIterable<string>
}

type CompletionOptions = {
  maxTokens?:   number   // Default: 1024
  temperature?: number   // Default: 0.1 (faktische Antworten)
  timeoutMs?:   number   // aus ConfigService
}
\end{lstlisting}

Modellname und Timeout werden als Config-Parameter gespeichert (\ra ADR-003):
\code{llm\_model}, \code{llm\_max\_tokens}, \code{llm\_temperature},
\code{llm\_timeout\_ms}. Provider-Wechsel innerhalb desselben Anbieters
erfordert nur eine DB-Änderung.

\subsection{Kostenabschätzung MVP}

\begin{table}[htb]
\small\centering
\begin{tabular}{@{}L{5.5cm} L{4cm} L{4cm}@{}}
\toprule
\rowcolor{tabhead}
 & \textbf{Azure GPT-4o-mini} & \textbf{Claude 3.5 Haiku} \\
\midrule
Input (pro 1M Token)           & $\approx$\$0.15  & $\approx$\$0.80  \\
Output (pro 1M Token)          & $\approx$\$0.60  & $\approx$\$4.00  \\
Kosten/Anfrage ($\approx$3.5k Token) & $\approx$\$0.001 & $\approx$\$0.005 \\
Kosten bei 1.000 Anfragen/Monat & $\approx$\$1.--  & $\approx$\$5.--  \\
\bottomrule
\end{tabular}
\end{table}

\section{Begründung}

\paragraph{Warum kein Self-Hosted-Modell?}
GPU-Infrastruktur liegt ausserhalb des 480h-Budgets. Qualität von
Open-Source-Modellen bei deutschen Fachdomänen-Dokumenten ist noch
nicht auf Cloud-Niveau. Der LLM-Adapter lässt diesen Wechsel Post-MVP
zu, ohne Business-Logik zu ändern.

\paragraph{Warum nicht OpenAI direkt (ohne Azure)?}
OpenAI API bietet keine europäische Datenresidenz-Garantie.
Wenn Compliance irrelevant ist, wäre OpenAI direkt technisch
gleichwertig mit Pfad~B (Anthropic) und kostengünstiger als Pfad~A.

\section{Betrachtete Alternativen}

\subsection{Alternative~1 -- OpenAI API direkt}

\zurueckgestellt Kein EU-Datenresidenz-Angebot. Wenn Compliance keine Rolle
spielt, technisch gleichwertig mit Pfad~B und kostengünstiger als Pfad~A.

\subsection{Alternative~2 -- Self-Hosted (Ollama + Llama~3.1 70B)}

\abgelehnt GPU-Infrastruktur ausserhalb des Budgets. Qualität bei deutschen
Fachdokumenten ungesichert. Betriebsaufwand für 1-FTE-Maintenance zu hoch.

\subsection{Alternative~3 -- Mistral AI (EU-Provider)}

\zurueckgestellt Europäische Datenresidenz vorhanden. Mistral-Modelle bei
deutschen Domänentexten weniger evaluiert. Kleineres Ökosystem.
Viable Post-MVP-Alternative.

\section{Konsequenzen}

\paragraph{Positiv}
\begin{itemize}
  \item LLM-Adapter-Interface hält Provider vollständig aus Business-Logik heraus
  \item Modell ohne Deployment wechselbar (ConfigService, ADR-003)
  \item Streaming in beiden Pfaden verfügbar --- $p_{95} \leq 10\,\text{s}$ erreichbar
  \item Kostenvolumen bei Pilot-Nutzung minimal
\end{itemize}

\paragraph{Negativ / Risiken}
\begin{itemize}
  \item \textbf{Quota-Risiko} (Risiko~3): Azure-Quota-Anfragen können 2--4~Wochen
        dauern; Anthropic-Pläne schneller verfügbar. Quota-Anfrage sofort stellen.
  \item API-Key darf nie in Code, Logs oder DB erscheinen (\ra QA-03, M10)
  \item Cloud-LLMs: kein garantiertes $p_{95}$-SLA; eigener Timeout und Circuit
        Breaker (ADR-002) als Absicherung nötig
\end{itemize}

\section{Abhängigkeiten}

\begin{table}[htb]
\small\centering
\begin{tabularx}{\linewidth}{L{4cm} L{3cm} X}
\toprule
\rowcolor{tabhead}
\textbf{Abhängigkeit} & \textbf{Typ} & \textbf{Hinweis} \\
\midrule
ADR-002 (LLMAdapter-Interface) & Vorausgesetzt  & Provider hinter Interface gekapselt \\
ADR-003 (ConfigService)        & Vorausgesetzt  & Modellname, Timeout als Config-Parameter \\
ADR-005 (RAG-Stack)            & Nachgelagert   & Embedding-Provider-Wahl hängt von LLM-Provider ab \\
Risiko~3                       & Direkt         & Quota-Anfrage blockiert ggf.\ Sprint-1-Start \\
Legal/IT-Abteilung             & Entscheider    & Compliance-Anforderung bis Ende Sprint~0 \\
\bottomrule
\end{tabularx}
\end{table}

\section{Offen / Nächste Schritte}

\begin{itemize}[label=\todo]
  \item \textbf{Bis Ende Sprint~0:} Legal/IT bestätigt Datenresidenz-Anforderung
        \ra Pfad~A (Azure) oder Pfad~B (Anthropic)
  \item \textbf{Sofort nach Entscheid:} Quota-Anfrage beim gewählten Provider stellen
  \item Prompt-Template für RAG-Antworten mit Quellenangabe ausarbeiten und im Spike testen
  \item Rate-Limit-Werte des Providers dokumentieren \ra Circuit-Breaker-Konfiguration (ADR-002)
\end{itemize}

%%============================================================
\chapter{RAG-Stack: pgvector und multilinguales Embedding-Modell}
\label{chap:adr005}
%%============================================================

\begin{tcolorbox}[enhanced, colback=adrlightblue, colframe=adrblue,
  boxrule=0.9pt, arc=3pt, left=8pt, right=8pt, top=6pt, bottom=6pt]
\begin{tabular}{@{}L{3.2cm}L{11.5cm}@{}}
\textbf{Status}      & \statusProposed \\[3pt]
\textbf{Datum}       & 2026-05-27 \\[3pt]
\textbf{Autor}       & Christoph A. Amstutz, MD-PhD \\[3pt]
\textbf{Bezug}       & US-01, US-04, QA-06 (Testability), Risiko~1 \\[3pt]
\textbf{Abhängigkeit}& ADR-001 (ConfidenceScorer), ADR-002 (RAG-Pipeline), ADR-004 (LLM-Provider) \\
\end{tabular}
\end{tcolorbox}

\section{Kontext}

Der RAG-Stack ist die technische Grundlage des Kernversprechens: verlässliche,
quellenbelegte Antworten. Risiko~1 stuft die RAG-Qualität als existenzielles
Go/No-Go-Risiko ein. Der Stack umfasst drei eng gekoppelte Entscheidungen:
Vektor-Datenbank, Embedding-Modell und Chunking-Strategie.

\begin{table}[htb]
\small\centering
\begin{tabularx}{\linewidth}{L{5cm} L{2cm} X}
\toprule
\rowcolor{tabhead}
\textbf{Anforderung} & \textbf{Quelle} & \textbf{Implikation} \\
\midrule
Dokumente: PDF, DOCX, Markdown       & US-04 & Parser für 3 Formate nötig \\
$\leq$50 Seiten $\rightarrow$ 5~min  & US-04 & Asynchrones Chunking + Embedding \\
Quellenreferenz: Titel + Abschnitt   & US-01 & Chunks müssen Herkunftsmetadaten tragen \\
Klick \ra Dokument mit Highlight     & US-01 & Chunks müssen Offset/Seitenzahl speichern \\
Sprache: Deutsch (MVP)               & Req.  & Embedding-Modell muss Deutsch unterstützen \\
Out-of-Corpus ,,Weiss nicht`` $\geq 90\,\%$ & US-02 & Retrieval-Score als Konfidenz-Signal \\
\bottomrule
\end{tabularx}
\end{table}

\section{Entscheidung}

\subsection{Komponente~1 -- Vektor-Datenbank: pgvector}

Für das MVP wird \textbf{pgvector} als Vektor-Datenbank eingesetzt ---
als Extension der bereits vorhandenen PostgreSQL-Instanz.

\begin{lstlisting}[language=LearnSQL]
CREATE EXTENSION IF NOT EXISTS vector;

CREATE TABLE document_chunks (
  id           BIGSERIAL    PRIMARY KEY,
  document_id  UUID         NOT NULL REFERENCES documents(id) ON DELETE CASCADE,
  chunk_index  INTEGER      NOT NULL,
  content      TEXT         NOT NULL,
  embedding    vector(1024) NOT NULL,
  page_number  INTEGER,                  -- fuer PDF-Highlight (US-01)
  char_offset  INTEGER,
  token_count  INTEGER      NOT NULL,
  created_at   TIMESTAMPTZ  NOT NULL DEFAULT now()
);

-- HNSW-Index fuer schnelle Aehnlichkeitssuche
CREATE INDEX ON document_chunks
  USING hnsw (embedding vector_cosine_ops)
  WITH (m = 16, ef_construction = 64);
\end{lstlisting}

\begin{lstlisting}[language=LearnSQL]
-- Aehnlichkeitssuche mit Score und Bereichsisolation (US-04)
SELECT dc.content, dc.page_number, dc.char_offset,
       d.title AS document_title, d.uploaded_at,
       1 - (dc.embedding <=> $1::vector) AS cosine_similarity
FROM   document_chunks dc
JOIN   documents d ON d.id = dc.document_id
WHERE  d.area_id = $2
ORDER  BY dc.embedding <=> $1::vector
LIMIT  $3;
\end{lstlisting}

\subsection{Komponente~2 -- Embedding-Modell: multilingual-e5-large-instruct}

\begin{table}[htb]
\small\centering
\begin{tabular}{@{}ll@{}}
\toprule
\rowcolor{tabhead}
\textbf{Eigenschaft} & \textbf{Wert} \\
\midrule
Modell              & \code{intfloat/multilingual-e5-large-instruct} \\
Embedding-Dimension & 1\,024 \\
Max.\ Input-Token   & 512 \\
Sprachen            & 100+, inkl.\ Deutsch \\
Modellgrösse        & $\approx$560~MB \\
CPU-Inferenz        & $\approx$150--300~ms / 512~Token \\
Lizenz              & MIT \\
\bottomrule
\end{tabular}
\end{table}

\begin{lstlisting}[language=Python]
class EmbeddingAdapter:
    def embed_query(self, query: str) -> list[float]:
        # E5-Modelle: unterschiedliche Prefixes fuer Query vs. Passage
        return self.model.encode(f"query: {query}").tolist()

    def embed_document(self, text: str) -> list[float]:
        return self.model.encode(f"passage: {text}").tolist()
\end{lstlisting}

\subsection{Komponente~3 -- Chunking-Strategie}

\begin{table}[htb]
\small\centering
\begin{tabular}{@{}ll@{}}
\toprule
\rowcolor{tabhead}
\textbf{Parameter} & \textbf{Wert} \\
\midrule
Strategie       & Rekursiver Splitter (Hierarchie: \textbackslash n\textbackslash n $\to$ \textbackslash n $\to$ Satz $\to$ Wort) \\
Chunk-Grösse    & 400~Token \\
Overlap         & 50~Token \\
Mindestgrösse   & 50~Token (filtert Artefakte) \\
\bottomrule
\end{tabular}
\end{table}

\subsection{Gesamt-Architektur}

\begin{lstlisting}[style=diagram]
Dokument-Upload (US-04)
        |
        v
+------------------------------------+
|  DocumentParser                    |
|  PDF  -> pdfplumber / pypdf         |
|  DOCX -> python-docx                |
|  MD   -> direkt                    |
+----------------+-------------------+
                 |
                 v
+------------------------------------+
|  ChunkingService                   |
|  Rekursiver Splitter               |
|  400 Token / 50 Overlap            |
+----------------+-------------------+
                 |
                 v
+------------------------------------+
|  EmbeddingAdapter                  |
|  multilingual-e5-large-instruct    |
+----------------+-------------------+
                 |
                 v
+------------------------------------+
|  VectorStore (pgvector)            |
|  INSERT document_chunks            |
|  HNSW-Index automatisch aktuell    |
+------------------------------------+


Suchanfrage (US-01)
        |
        v EmbeddingAdapter.embed_query()
        v
+------------------------------------+
|  Retriever                         |
|  SELECT ... ORDER BY cosine dist   |
|  WHERE area_id = {bereich}         |
|  LIMIT top_k (Default: 5)          |
+------------------------------------+
        |
        v --> RAG-Pipeline (ADR-002)
\end{lstlisting}

\subsection{Retriever-Interface}

\begin{lstlisting}[language=TypeScript]
interface Retriever {
  search(query: string, areaId: string, topK?: number): Promise<RetrievedChunk[]>
}

type RetrievedChunk = {
  content:       string
  similarity:    number        // 0.0 -- 1.0
  documentTitle: string
  uploadedAt:    Date
  pageNumber:    number | null
  charOffset:    number
  documentId:    string
}

class InMemoryRetriever implements Retriever {
  constructor(private fixtures: RetrievedChunk[]) {}
  async search(): Promise<RetrievedChunk[]> { return this.fixtures }
}
\end{lstlisting}

\section{Begründung}

\paragraph{Warum pgvector statt dedizierter Vektor-DB (Qdrant, Chroma)?}
pgvector nutzt die bereits vorhandene PostgreSQL-Instanz --- keine neue
Infrastruktur, kein zusätzliches Betriebswissen. Der HNSW-Index ist für
500~Dokumente und 20 concurrent users performant genug. Der
Retriever-Interface-Layer erlaubt den Wechsel zu Qdrant transparent.

\paragraph{Warum lokales Embedding-Modell statt API?}
(1)~Kein zweiter externer Service = keine zweite Quota-Abhängigkeit (Risiko~3).
(2)~Keine Embedding-Kosten beim Re-Indexieren.
(3)~Dokumenteninhalt verlässt das System nicht --- konsistente Datenresidenz
wenn ADR-004 auf Azure OpenAI zeigt.

\paragraph{Warum multilingual-e5-large und nicht kleiner?}
\code{multilingual-e5-large} übertrifft \code{e5-small} bei deutschen
Domänentexten deutlich (MTEB-Benchmark). 560~MB Modellgrösse sind bei
Server-Deployment kein Problem. \code{multilingual-e5-base} ist Fallback
wenn CPU-Inferenzzeit zu hoch.

\section{Betrachtete Alternativen}

\subsection{Alternative~1 -- Qdrant (dediziert, self-hosted)}

\zurueckgestellt Für MVP-Corpus (500~Dokumente, 1~Bereich) Overkill.
Neuer Infrastruktur-Container, Betriebsaufwand. Upgrade-Pfad via Interface.

\subsection{Alternative~2 -- Chroma (lokal)}

\abgelehnt Kein SQL-Join möglich (Bereichsisolation). Transaktionssicherheit
schlechter als PostgreSQL. Filesystem-Persistenz problematisch bei Backups.

\subsection{Alternative~3 -- OpenAI text-embedding-3-small (API)}

\zurueckgestellt Zweite externe API-Abhängigkeit (Risiko~3).
Datenschutzinkonsistenz wenn LLM-Provider Azure ist, Embeddings aber
zu OpenAI~US gehen. Viable wenn lokale Compute-Infrastruktur Problem ist.

\section{Konsequenzen}

\paragraph{Positiv}
\begin{itemize}
  \item Keine neue Infrastrukturkomponente für MVP
  \item Bereichsisolation nativ in SQL (\code{WHERE area\_id = ?})
  \item Lokales Embedding: keine API-Kosten, keine externen Daten beim Upload
  \item \code{InMemoryRetriever} macht alle Tests ohne DB und Embedding-Modell ausführbar
\end{itemize}

\paragraph{Negativ / Risiken}
\begin{itemize}
  \item pgvector HNSW: bei $> 1\,\text{M}$ Chunks degradiert Performance.
        Mitigation: Upgrade auf Qdrant via Interface
  \item Lokales Modell: $\approx$560~MB auf dem Server; beim Start laden
        (kein Cold-Start pro Request)
  \item \textbf{Go/No-Go-Risiko (Risiko~1):} Retrieval-Qualität auf echten
        Dokumenten muss im Spike vor Sprint~1 validiert werden
\end{itemize}

\section{Abhängigkeiten}

\begin{table}[htb]
\small\centering
\begin{tabularx}{\linewidth}{L{4cm} L{3cm} X}
\toprule
\rowcolor{tabhead}
\textbf{Abhängigkeit} & \textbf{Typ} & \textbf{Hinweis} \\
\midrule
ADR-002 (RAG-Pipeline)    & Konsument     & \code{Retriever.search()} ist erster Pipeline-Schritt \\
ADR-001 (ConfidenceScorer)& Konsument     & \code{SemanticSimilarityScorer} nutzt denselben \code{EmbeddingAdapter} \\
ADR-004 (LLM-Provider)    & Koordiniert   & Datenresidenz für Embeddings muss konsistent sein \\
ADR-003 (ConfigService)   & Konsument     & \code{top\_k}, Ähnlichkeitsschwelle als Config-Parameter \\
Risiko~1                  & Go/No-Go      & Spike mit Evaluierungs-Harness vor Sprint~1 \\
\bottomrule
\end{tabularx}
\end{table}

\section{Offen / Nächste Schritte}

\begin{itemize}[label=\todo]
  \item \textbf{Sprint~0 Spike:} multilingual-e5-large mit echten Pilotdokumenten
        und Evaluierungs-Dataset testen; Source-Hit-Rate, Out-of-Corpus-Rate,
        $p_{95}$-Latenz messen
  \item Chunk-Grösse variieren (256 / 400 / 512~Token) und optimalen Wert messen
  \item PDF-Tabellenextraktion evaluieren: pdfplumber vs.\ pymupdf für Layoutqualität
  \item Re-Indexierung bei Dokumenten-Update (gleicher Dateiname, US-04):
        alte Chunks löschen, neue insertieren --- in einer Transaktion
  \item \code{top\_k} und Ähnlichkeitsschwellenwert in \code{config\_params}
        eintragen (\ra ADR-003)
\end{itemize}

%%============================================================
\chapter{Authentifizierungs-Strategie: JWT mit HTTP-only Cookie und SSO-Ready-Middleware}
\label{chap:adr006}
%%============================================================

\begin{tcolorbox}[enhanced, colback=adrlightblue, colframe=adrblue,
  boxrule=0.9pt, arc=3pt, left=8pt, right=8pt, top=6pt, bottom=6pt]
\begin{tabular}{@{}L{3.2cm}L{11.5cm}@{}}
\textbf{Status}      & \statusProposed \\[3pt]
\textbf{Datum}       & 2026-05-27 \\[3pt]
\textbf{Autor}       & Christoph A. Amstutz, MD-PhD \\[3pt]
\textbf{Bezug}       & US-05, QA-03 (Security), QA-05 (Maintainability) \\[3pt]
\textbf{Abhängigkeit}& ADR-007 (Frontend-Framework) \\
\end{tabular}
\end{tcolorbox}

\section{Kontext}

US-05 definiert die MVP-Authentifizierungsanforderungen:
\begin{itemize}
  \item Username/Passwort (lokal), Accounts per DB-Script
  \item Keine Self-Service-Registrierung
  \item Drei Rollen: \code{Lernende} / \code{Bereichsverantwortlicher} / \code{Admin}
  \item Nicht authentifizierte Requests \ra Redirect auf Login-Seite
  \item \textbf{Post-MVP:} SSO via Azure AD / SAML~2.0 oder OIDC
\end{itemize}

Die Post-MVP-SSO-Anforderung ist der entscheidende Architektur-Treiber: die
Auth-Strategie muss so gebaut werden, dass der Wechsel zu SSO kein Refactoring
der Business-Logik, der RBAC-Middleware oder der geschützten Endpunkte erfordert.

\section{Entscheidung}

\textbf{JWT-basierte Authentifizierung} mit kurzlebigem Access-Token (15~min)
und langlebigem Refresh-Token (7~Tage), beide als \textbf{HTTP-only Cookies}.
Die gesamte Auth-Logik lebt in einer \code{AuthMiddleware}-Schicht hinter einem
Interface. RBAC ist in einer separaten \code{RBACMiddleware} gekapselt.

\subsection{Token-Strategie}

\begin{table}[htb]
\small\centering
\begin{tabular}{@{}L{2.8cm} L{2.5cm} L{2cm} L{3cm} L{3.5cm}@{}}
\toprule
\rowcolor{tabhead}
\textbf{Token} & \textbf{Typ} & \textbf{Dauer} & \textbf{Transport} & \textbf{Inhalt} \\
\midrule
Access-Token  & JWT (HS256) & 15 min  & HTTP-only Cookie & \code{userId}, \code{role}, \code{areaId}, \code{exp} \\
Refresh-Token & Opaque (UUID) & 7 Tage & HTTP-only Cookie & in DB gespeichert \\
\bottomrule
\end{tabular}
\end{table}

\begin{lstlisting}[language=TypeScript]
type AccessTokenPayload = {
  sub:    string   // userId (UUID)
  role:   Role     // 'learner' | 'area_manager' | 'admin'
  areaId: string
  exp:    number   // Unix Timestamp (now + 15 min)
}
type Role = 'learner' | 'area_manager' | 'admin'
\end{lstlisting}

\subsection{Request-Flow}

\begin{lstlisting}[style=diagram]
Request
    |
    v
+-----------------------------------------------------+
|  AuthMiddleware (Interface)                         |
|                                                     |
|  MVP: JWT aus HTTP-only Cookie validieren           |
|       -> request.user = { id, role, areaId }        |
|                                                     |
|  Post-MVP SSO: OIDC/SAML-Token validieren           |
|       -> request.user setzen (gleiche Schnittstelle)|
+----------------------+------------------------------+
                       | request.user gesetzt
                       v
+-----------------------------------------------------+
|  RBACMiddleware (unveraendert beim SSO-Wechsel)     |
|  pruefte: request.user.role in allowed_roles        |
|  -> 403 Forbidden wenn nicht berechtigt             |
+----------------------+------------------------------+
                       |
                       v
               Business Handler
\end{lstlisting}

\subsection{AuthMiddleware-Interface}

\begin{lstlisting}[language=TypeScript]
interface AuthMiddleware {
  authenticate(req: Request, res: Response, next: NextFunction): Promise<void>
}

class LocalJWTAuthMiddleware  implements AuthMiddleware { ... }  // MVP
class OIDCAuthMiddleware      implements AuthMiddleware { ... }  // Post-MVP

class MockAuthMiddleware implements AuthMiddleware {
  constructor(private user: AuthUser) {}
  async authenticate(req: any, res: any, next: NextFunction) {
    req.user = this.user; next()
  }
}
\end{lstlisting}

\subsection{RBAC-Deklaration pro Route}

\begin{lstlisting}[language=TypeScript]
router.post('/api/documents/upload',
  rbac(['area_manager', 'admin']), uploadHandler)

router.get('/api/admin/config',
  rbac(['admin']), configHandler)

function rbac(allowedRoles: Role[]): RequestHandler {
  return (req, res, next) => {
    if (!allowedRoles.includes((req as any).user?.role))
      return res.status(403).json({ error: 'Forbidden' })
    next()
  }
}
\end{lstlisting}

\subsection{Passwort-Hashing (Argon2id)}

\begin{lstlisting}[language=TypeScript]
import argon2 from 'argon2'

const OPTIONS = { type: argon2.argon2id, memoryCost: 65536,
                  timeCost: 3, parallelism: 4 }

async function hashPassword(plain: string):   Promise<string>  {
  return argon2.hash(plain, OPTIONS)
}
async function verifyPassword(hash: string, plain: string): Promise<boolean> {
  return argon2.verify(hash, plain)
}
\end{lstlisting}

\subsection{Datenbankschema}

\begin{lstlisting}[language=LearnSQL]
CREATE TABLE users (
  id            UUID    PRIMARY KEY DEFAULT gen_random_uuid(),
  username      TEXT    NOT NULL UNIQUE,
  password_hash TEXT    NOT NULL,
  role          TEXT    NOT NULL CHECK (role IN ('learner','area_manager','admin')),
  area_id       TEXT    NOT NULL,
  created_at    TIMESTAMPTZ NOT NULL DEFAULT now(),
  is_active     BOOLEAN NOT NULL DEFAULT true
);

CREATE TABLE refresh_tokens (
  id         UUID    PRIMARY KEY DEFAULT gen_random_uuid(),
  user_id    UUID    NOT NULL REFERENCES users(id) ON DELETE CASCADE,
  token_hash TEXT    NOT NULL UNIQUE,   -- SHA-256 des Tokens
  expires_at TIMESTAMPTZ NOT NULL,
  created_at TIMESTAMPTZ NOT NULL DEFAULT now(),
  revoked_at TIMESTAMPTZ              -- NULL = aktiv
);
\end{lstlisting}

\section{Begründung}

\paragraph{Warum JWT + HTTP-only Cookie statt Server-Side Session?}
JWT-Access-Tokens sind zustandslos --- Validierung ohne DB in $< 1\,\text{ms}$.
Der Refresh-Token in der DB ermöglicht trotzdem Server-seitige Invalidierung
bei Logout oder Sicherheitsvorfällen.

\paragraph{Warum HTTP-only Cookie statt Authorization-Header?}
HTTP-only Cookies sind für JavaScript nicht lesbar --- XSS-Angriffe können den
Token nicht stehlen. Für ein internes Tool ohne öffentliche API ist dieser
Trade-off optimal.

\paragraph{Warum Argon2id statt Bcrypt?}
Argon2id ist der aktuelle OWASP-Standard (2024) und resistent gegen
GPU-basierte Brute-Force-Angriffe durch Memory-Hardness. Kosten: eine
npm-Library.

\paragraph{Warum Refresh-Token-Rotation?}
Bei Rotation wird nach jedem Refresh der alte Token invalidiert und ein neuer
ausgestellt. Ein gestohlener Refresh-Token wird beim nächsten legitimen Refresh
erkannt (Collision \ra beide Sessions invalidieren).

\paragraph{Warum \code{AuthMiddleware} als Interface?}
Die Post-MVP-SSO-Anforderung ist kein Nice-to-Have. Ohne Interface zieht sich
SAML/OIDC-Logik durch alle geschützten Endpunkte. Das Interface kapselt die
Auth-Logik an genau einem Ort.

\section{Betrachtete Alternativen}

\subsection{Alternative~1 -- Auth.js (NextAuth.js)}

\zurueckgestellt Auth.js bietet native Next.js-Integration mit eingebautem
OIDC/SAML-Support. Wenn ADR-007 Next.js wählt, sollte Auth.js als primäre
Auth-Integration evaluiert werden --- weniger Eigenimplementierung,
dafür weniger Kontrolle über Token-Format.
\textbf{Empfehlung:} Auth.js bevorzugen wenn Next.js als Framework gewählt.

\subsection{Alternative~2 -- Server-Side Session (express-session)}

\abgelehnt DB-Lookup bei jedem Request. Horizontal schwerer skalierbar.
SSO-Readiness erfordert dasselbe Interface-Konzept.

\subsection{Alternative~3 -- Keycloak (Self-Hosted IdP)}

\abgelehnt Separater Server, erheblicher Setup- und Betriebsaufwand für
480h-MVP-Budget. Sinnvoll wenn das Unternehmen bereits einen zentralen IdP betreibt.

\section{Konsequenzen}

\paragraph{Positiv}
\begin{itemize}
  \item SSO-Wechsel (Post-MVP) erfordert ausschliesslich neue
        \code{AuthMiddleware}-Implementierung
  \item \code{MockAuthMiddleware} macht alle geschützten Endpunkte ohne
        Login-Flow testbar
  \item HTTP-only Cookies verhindern XSS-Token-Diebstahl
  \item Argon2id: state-of-the-art Passwortschutz
  \item Token-Rotation erkennt gestohlene Refresh-Tokens
\end{itemize}

\paragraph{Negativ / Risiken}
\begin{itemize}
  \item JWT-Invalidierung vor Ablauf nicht möglich ohne DB-Lookup.
        Mitigation: Refresh-Token revoken \ra nächste Rotation schlägt fehl
  \item \code{HS256} (symmetrisch): JWT-Secret muss sicher im Secrets Manager
        liegen, nie in Code. Bei Multi-Instance wäre \code{RS256} besser ---
        für MVP-Single-Instance akzeptabel
  \item Argon2id-Hashing dauert $\approx$100--300~ms (gewollt als
        Brute-Force-Schutz) --- Login-Latenz entsprechend
\end{itemize}

\section{Abhängigkeiten}

\begin{table}[htb]
\small\centering
\begin{tabularx}{\linewidth}{L{4cm} L{3cm} X}
\toprule
\rowcolor{tabhead}
\textbf{Abhängigkeit} & \textbf{Typ} & \textbf{Hinweis} \\
\midrule
ADR-007 (Frontend-Framework)  & Koordiniert   & Wenn Next.js \ra Auth.js-Alternative bevorzugen \\
US-05                         & Direkt        & Alle ACs von US-05 werden abgedeckt \\
QA-03 Security (M8, M9)       & Validierung   & RBAC server-side setzt diese Auth-Schicht voraus \\
ADR-003 (ConfigService)       & Optional      & JWT-Secret-Lebensdauer als Config-Parameter möglich \\
\bottomrule
\end{tabularx}
\end{table}

\section{Offen / Nächste Schritte}

\begin{itemize}[label=\todo]
  \item ADR-007 abschliessen: wenn Next.js \ra Auth.js-Alternative evaluieren
  \item JWT-Secret-Management definieren: Vault, Docker Secret oder ENV-Variable?
  \item DB-Script für erste Admin- und Testaccounts schreiben (Argon2id-Hashes)
  \item SSO-Provider-Typ (SAML~2.0 vs.\ OIDC) für Post-MVP-Planung klären
\end{itemize}

%%============================================================
\chapter{Frontend-Framework: Next.js mit TypeScript und App Router}
\label{chap:adr007}
%%============================================================

\begin{tcolorbox}[enhanced, colback=adrlightblue, colframe=adrblue,
  boxrule=0.9pt, arc=3pt, left=8pt, right=8pt, top=6pt, bottom=6pt]
\begin{tabular}{@{}L{3.2cm}L{11.5cm}@{}}
\textbf{Status}      & \statusProposed \\[3pt]
\textbf{Datum}       & 2026-05-27 \\[3pt]
\textbf{Autor}       & Christoph A. Amstutz, MD-PhD \\[3pt]
\textbf{Bezug}       & US-01, US-03, US-07--09, US-11, QA-05 (Maintainability) \\[3pt]
\textbf{Abhängigkeit}& ADR-002 (Discriminated Union / TypeScript), ADR-006 (Auth-Strategie) \\
\end{tabular}
\end{tcolorbox}

\section{Kontext}

LearnFlow benötigt vier funktional unterschiedliche UI-Bereiche:

\begin{table}[htb]
\small\centering
\begin{tabularx}{\linewidth}{L{3.5cm} X L{3.5cm}}
\toprule
\rowcolor{tabhead}
\textbf{Bereich} & \textbf{Hauptfunktionen} & \textbf{Besonderheiten} \\
\midrule
Chat-Interface (Lara) & Frage stellen, Antwort + Quellenlinks, Feedback, Folgefragen & LLM-Streaming \\
Lern-Interface (Lara) & Quiz absolvieren, Ergebnis & --- \\
Dashboard (Stefan)    & Dokumente verwalten, Quiz-Freigabe, Wissenslücken & Upload, Review \\
Admin-Seite           & Schwellenwerte, Audit-Log & Formular-getrieben \\
\bottomrule
\end{tabularx}
\end{table}

Aus ADR-002 ergibt sich eine technische Vorentscheidung: der
Discriminated-Union-Exhaustiveness-Check ist nativ und vollständig
typsicher nur in \textbf{TypeScript} umsetzbar. Das Framework muss
TypeScript unterstützen.

\textbf{Scope: Desktop-Browser, kein Mobile als MVP-Anforderung.}

\section{Entscheidung}

\textbf{Next.js~15 (App Router) mit TypeScript} als Full-Stack-Framework.
Next.js übernimmt sowohl das Frontend (React-Komponenten) als auch
das Backend (API-Routes) --- kein separater API-Server für das MVP nötig.

\subsection{Architektur-Übersicht}

\begin{lstlisting}[style=diagram]
Browser
    | HTTPS
    v
+--------------------------------------------------+
|  Next.js 15 (App Router, TypeScript)             |
|                                                  |
|  app/                                            |
|    (auth)/login/          -> Login (public)      |
|    (learner)/chat/        -> Chat (Lara)         |
|    (learner)/quiz/        -> Quiz (Lara)         |
|    (manager)/documents/   -> Dokumente (Stefan)  |
|    (manager)/quiz-review/ -> Quiz-Freigabe       |
|    (manager)/gaps/        -> Wissensluecken      |
|    (admin)/config/        -> Admin-Seite         |
|                                                  |
|  app/api/                 -> API-Routes          |
|    auth/{login,refresh,logout}                   |
|    questions/             -> RAG-Pipeline        |
|    documents/             -> Corpus-Verwaltung   |
|    quiz/, admin/config/                          |
+----------------------+---------------------------+
                       | HTTP (interner Aufruf)
                       v
           +-------------------------+
           |  Python Backend         |
           |  RAG-Engine, Embedding  |
           |  Chunking, pgvector     |
           +-------------------------+
\end{lstlisting}

Die Python-Backend-Service-Grenze trennt das ML-Ökosystem
(\code{sentence-transformers}, \code{pdfplumber}, \code{python-docx})
vom TypeScript-Frontend. Das ist die einzige Service-Grenze im MVP.

\subsection{TypeScript als Pflicht (Verbindung zu ADR-002)}

\begin{lstlisting}[language=TypeScript]
// Compiler-Fehler wenn ein RAGResult-Variant nicht behandelt wird
function renderResult(result: RAGResult): React.ReactNode {
  switch (result.kind) {
    case 'answer':      return result.label === 'limited'
                          ? <LimitedAnswerCard {...result} />
                          : <AnswerCard {...result} />
    case 'dont_know':   return <DontKnowCard hint={result.hint} />
    case 'no_source':
    case 'suppressed':  return <NoSourceCard />
    case 'unavailable': return <ErrorCard message={result.message} />
    case 'timeout':     return <ErrorCard message="Zeitüberschreitung." />
    // Fehlender Case => TypeScript-Compilerfehler
  }
}
\end{lstlisting}

\subsection{LLM-Streaming im Chat-Interface (App Router)}

\begin{lstlisting}[language=TypeScript]
// app/api/questions/route.ts
export async function POST(req: Request) {
  const { question, history } = await req.json()

  const stream = new ReadableStream({
    async start(controller) {
      const ragResult = await runRAGPipeline(question, history)
      if (ragResult.kind !== 'answer') {
        controller.enqueue(JSON.stringify(ragResult))
        controller.close(); return
      }
      for await (const chunk of llmAdapter.stream(ragResult.prompt))
        controller.enqueue(chunk)
      controller.close()
    }
  })
  return new Response(stream, { headers: { 'Content-Type': 'text/event-stream' }})
}
\end{lstlisting}

\subsection{Auth.js-Integration für Post-MVP SSO}

\begin{lstlisting}[language=TypeScript]
// auth.ts (wenn Auth.js gewaehlt -- Alternative 1 aus ADR-006)
import NextAuth from 'next-auth'
import Credentials from 'next-auth/providers/credentials'

export const { handlers, auth, signIn, signOut } = NextAuth({
  providers: [
    Credentials({
      async authorize({ username, password }) {
        return await verifyUser(username as string, password as string) ?? null
      }
    })
    // Post-MVP: SAML/OIDC-Provider hier ergaenzen -- kein Code unterhalb aendert sich
  ],
  callbacks: {
    jwt({ token, user }) {
      if (user) { token.role = user.role; token.areaId = user.areaId }
      return token
    }
  }
})
\end{lstlisting}

\subsection{Verzeichnisstruktur}

\begin{lstlisting}[style=diagram]
learnflow/
  app/
    (auth)/login/page.tsx
    (learner)/chat/page.tsx
    (learner)/quiz/page.tsx
    (manager)/documents/
    (manager)/quiz-review/
    (admin)/config/
    api/{auth,questions,documents,quiz,admin}/
  components/
    chat/{AnswerCard,LimitedAnswerCard,DontKnowCard,NoSourceCard,ErrorCard}.tsx
    shared/
  lib/
    rag/          -- RAGResult-Typen, Pipeline-Client
    auth/         -- AuthMiddleware-Interface + Implementierung
    config/       -- ConfigService-Client
  types/rag.ts    -- RAGResult Discriminated Union (ADR-002)
  middleware.ts   -- Auth-Schutz aller Routen
\end{lstlisting}

\section{Begründung}

\paragraph{Warum Next.js statt React SPA + Express?}
Zwei separate Services zu deployen und zu betreiben erhöht den Aufwand für
ein 1-FTE-Maintenance-Team erheblich. Next.js (ein Deployment, kein CORS,
ein Dev-Server) ist die ökonomischere Wahl. API-Routes sind vollwertige
Node.js-Handler --- kein Funktionalitätsverlust.

\paragraph{Warum App Router statt Pages Router?}
App Router bietet native Streaming-Unterstützung via React~18 Suspense und
\code{ReadableStream}. Das ist für die LLM-Streaming-Anforderung
($p_{95} \leq 10\,\text{s}$, US-09) der entscheidende Vorteil.

\paragraph{Warum nicht SvelteKit?}
Svelte ist schneller und leichtgewichtiger. Der Hauptgrund dagegen:
der Discriminated-Union-Exhaustiveness-Check (ADR-002) ist in Svelte-Templates
nicht compiler-erzwingbar. TypeScript-Sicherheit ist hier keine optionale
Verbesserung, sondern eine Kernanforderung.

\paragraph{Warum Python für Backend-Services?}
Das ML-Ökosystem für Embedding (\code{sentence-transformers}), PDF-Parsing
(\code{pdfplumber}) und zukünftige Scoring-Modelle ist in Python deutlich
reifer. Eine einzelne Service-Grenze hält diese Stärke nutzbar ohne die
Frontend-Entwicklung auf Python zu zwingen.

\section{Betrachtete Alternativen}

\subsection{Alternative~1 -- SvelteKit + TypeScript}

\abgelehnt Exhaustiveness-Check für Discriminated Unions in Svelte-Templates
nicht erzwingbar. Kleineres Ökosystem für Auth (Post-MVP SSO).

\subsection{Alternative~2 -- React SPA (Vite) + Express.js}

\abgelehnt Zwei separate Services. CORS-Konfiguration. Kein natives Streaming
ohne manuelle SSE-Implementierung. Kein Mehrwert gegenüber Next.js.

\subsection{Alternative~3 -- Vue~3 + Nuxt}

\abgelehnt Kein technischer Vorteil gegenüber Next.js. Auth.js-Integration
für Nuxt schwächer als für Next.js.

\section{Konsequenzen}

\paragraph{Positiv}
\begin{itemize}
  \item Ein einziges Deployment für Frontend und BFF-API
  \item TypeScript end-to-end: Discriminated Unions aus ADR-002 voll typsicher
  \item Natives Streaming für LLM-Antworten (App Router + React~18)
  \item Auth.js deckt Post-MVP SSO ohne Refactoring
  \item Route-Groups (\code{(admin)}, \code{(learner)}) isolieren RBAC
\end{itemize}

\paragraph{Negativ / Risiken}
\begin{itemize}
  \item App Router jünger als Pages Router --- mehr API-Instabilität in edge cases
  \item Python-Backend als zweiter Service: eine Service-Grenze; Mitigation:
        Docker~Compose für lokale Entwicklung
  \item Server~Components vs.\ Client~Components: erhöhter Erklärungsaufwand
        beim Onboarding
\end{itemize}

\section{Abhängigkeiten}

\begin{table}[htb]
\small\centering
\begin{tabularx}{\linewidth}{L{4cm} L{3cm} X}
\toprule
\rowcolor{tabhead}
\textbf{Abhängigkeit} & \textbf{Typ} & \textbf{Hinweis} \\
\midrule
ADR-002 (Discriminated Union) & Voraussetzung & TypeScript ist Pflicht --- Next.js erfüllt das \\
ADR-006 (Auth-Strategie)      & Koordiniert   & Auth.js als bevorzugte MVP-Auth \\
ADR-005 (RAG-Stack / Python)  & Nachgelagert  & Next.js API-Routes kommunizieren via HTTP mit Python-Backend \\
US-09 (Folgefragen)           & Feature       & Konversationshistorie (max.\ 3 Paare) im Client-State \\
\bottomrule
\end{tabularx}
\end{table}

\section{Offen / Nächste Schritte}

\begin{itemize}[label=\todo]
  \item ADR-006 aktualisieren: Auth.js als primäre Auth-Integration bestätigen
  \item Python-Backend-Interface definieren: OpenAPI-Spec oder TypeScript-Types
        für den HTTP-Kontrakt
  \item Docker-Compose-Setup: Next.js + Python-Backend + PostgreSQL (mit pgvector)
  \item Styling-Strategie entscheiden: shadcn/ui oder Tailwind CSS direkt
  \item Streaming-Implementierung prototypisch testen ---
        $p_{95}$-Messung mit Mock-LLM und realem Netzwerk
\end{itemize}

%%============================================================
\end{document}
